\documentclass{tufte-handout}

%\geometry{showframe}
%\geometry{showframe}% for debugging purposes -- displays the margins

\usepackage{amsmath}
\usepackage{amssymb}
\usepackage{cleveref}
\crefname{Proposition}{Proposition}{Propositions}
\Crefname{Proposition}{Proposition}{Propositions}
\crefname{Theorem}{Theorem}{Theorems}
\Crefname{Theorem}{Theorem}{Theorems}
\crefname{Definition}{Definition}{Definitions}
\Crefname{Definition}{Definition}{Definitions}
\crefname{Corollary}{Corollary}{Corollaries}
\Crefname{Corollary}{Corollary}{Corollaries}
\crefname{Lemma}{Lemma}{Lemmas}
\Crefname{Lemma}{Lemma}{Lemmas}
\crefname{Example}{Example}{Examples}
\Crefname{Example}{Example}{Examples}
\usepackage{amsthm}

% Set up the images/graphics package
\usepackage{graphicx}
\setkeys{Gin}{width=\linewidth,totalheight=\textheight,keepaspectratio}
\graphicspath{{graphics/}}

% The following package makes prettier tables.  We're all about the bling!
\usepackage{booktabs}

% The units package provides nice, non-stacked fractions and better spacing
% for units.
\usepackage{units}

% The fancyvrb package lets us customize the formatting of verbatim
% environments.  We use a slightly smaller font.
\usepackage{fancyvrb}
\fvset{fontsize=\normalsize}

% Small sections of multiple columns
\usepackage{multicol}

% Provides paragraphs of dummy text
\usepackage{lipsum}

% Defines colors
\usepackage{xcolor}
\definecolor{blue}{cmyk}{0.63, 0.37, 0, 0.57}     
\definecolor{ltblue}{RGB}{78,150,179}

% These commands are used to pretty-print LaTeX commands
\newcommand{\doccmd}[1]{\texttt{\textbackslash#1}}% command name -- adds backslash automatically
\newcommand{\docopt}[1]{\ensuremath{\langle}\textrm{\textit{#1}}\ensuremath{\rangle}}% optional command argument
\newcommand{\docarg}[1]{\textrm{\textit{#1}}}% (required) command argument
\newenvironment{docspec}{\begin{quote}\noindent}{\end{quote}}% command specification environment
\newcommand{\docenv}[1]{\textsf{#1}}% environment name
\newcommand{\docpkg}[1]{\texttt{#1}}% package name
\newcommand{\doccls}[1]{\texttt{#1}}% document class name
\newcommand{\docclsopt}[1]{\texttt{#1}}% document class option name

% Package for title style
\usepackage{sectsty}
\usepackage[utf8]{inputenc}

% Sets section number style
\setcounter{secnumdepth}{3} % uncomment this, if desired
\renewcommand\thesection{\color{white}\arabic{section}}

% Sets title style
\makeatletter
  \renewcommand{\paragraph}{\@startsection{paragraph}%
    {4}{\z@}{-1ex \@plus -1ex \@minus -.3ex}%
    {0.5ex \@plus .2ex}{\normalfont\normalsize\bfseries}}
\makeatother

\makeatletter
  \renewcommand{\subsection}{\@startsection{subsection}%
    {3}{-1.8em}{-3ex \@plus -1ex \@minus -.2ex}%
    {1.5ex \@plus .2ex}
    {\hspace*{-5.5em}\fcolorbox{ltblue}{ltblue}{\parbox[c][1.0ex][b]{4em}{\phantom{space}}}
    \normalfont\large\itshape\color{ltblue}}}
\makeatother

\makeatletter
  \renewcommand{\section}{\@startsection{section}%
    {3}{-1.01em}{-3ex \@plus -1ex \@minus -.2ex}%
    {1.5ex \@plus .2ex}
    {\hspace*{-5.5em}\fcolorbox{blue}{blue}{\parbox[c][1.0ex][b]{4em}{\phantom{space}}}
    \normalfont\Large\itshape\color{blue}}}
\makeatother

% Sets theorem style
\usepackage{thmtools}
\usepackage{transparent}
\definecolor{theb}{rgb}{0.67, 0.80, 0.91}

\declaretheorem[shaded={bgcolor=Lavender,
    textwidth=30em}]{Definition} % Colorbox-styled Theorem 
\declaretheorem[shaded={bgcolor=Thistle,
    textwidth=30em}]{Theorem} % Colorbox-styled Theorem 
\declaretheorem[shaded={bgcolor=Thistle,
    textwidth=30em}]{Corollary} % Colorbox-styled Theorem
\declaretheorem[shaded={bgcolor=PeachPuff,
    textwidth=30em}]{Proposition} % Colorbox-styled Proposition
\declaretheorem[shaded={bgcolor=Thistle,
    textwidth=30em}]{Lemma} % Colorbox-styled Lemma
\declaretheorem[shaded={rulecolor=Lavender,
    rulewidth=2pt, bgcolor={rgb}{1,1,1}}]{Example-1} % Colorbounded-styled Theorem
\declaretheorem[thmbox=L]{boxtheorem L} % Theorem box L-size
\declaretheorem[thmbox=M]{Example} % Theorem box M-size
\declaretheorem[thmbox=S]{Formula} % Theorem box S-size
\declaretheorem[thmbox=S]{Statement} % Theorem box S-size

% set up pdf bookmark depth
\hypersetup{bookmarksdepth=3}

% ------------------------------------------------------------
\title{MATH563: Convex Optimization}
\author[Matthew He]{Matthew He}
  % if the \date{} command is left out, the current date will be used
% Beginning of the document
\begin{document}
\maketitle% this prints the handout title, author, and date

% abstract
\begin{abstract}
\noindent Abstract
\end{abstract}

% main text
\section{Preliminaries}
\begin{Definition}[Affine set]
    A set \( C \subseteq \mathbb{E} \) is called an
    affine set if for all \( x, y \in C \) and \( \lambda,\mu \in \mathbb{R} \), we have
    \begin{align}
    \lambda x + \mu y \in C.
    \end{align}
\end{Definition}

\begin{Definition}[Affine hull]
    The affine hull of a set \( Q \subseteq \mathbb{E} \)
    denoteded by \( \text{aff}(Q) \), is the intersection of all affine sets that 
    contain \( Q \). 
\end{Definition}
Affine hull is the smallest affine set that contains \( Q \).

\begin{Definition}[Relative interior and boundary]
    The relative interior of a set \( Q \subseteq \mathbb{E} \),
    denoted by \( \text{ri}(Q) \), is the interior of \( Q \) relative to its affine hull, i.e.
    \begin{align}
    \text{ri}(Q) = \{ x \in Q : \exists \delta > 0, B(x, \delta) \cap \text{aff}(Q) \subseteq Q \}.
    \end{align}
    The relative boundary of \( Q \) is defined as \( \text{rb}(Q) = \text{cl}(Q) \setminus \text{ri}(Q) \).
\end{Definition}
\marginnote{Relative interior is weaker than interior since it's easier to be satisfied, 
as \( B_\varepsilon \cap \text{aff} Q \subset B_{\varepsilon} \).
We have \( \text{int}(Q) \subseteq \text{ri}(Q) \) and \( \text{rb}(Q) \subseteq \text{bd}(Q) \).}

% \subsection{Differentiation in Euclidean space}
% \begin{Definition}[Second derivative]
%     If ll√
% \end{Definition}
\begin{Definition}[Epigraph]
    The epigraph of a function \( f:\mathbb{E}\to \overline{\mathbb{R}} \) is the set
    \begin{align}
    \text{epi}(f) = \{ (x, \alpha) \in \mathbb{E} \times \mathbb{R} : f(x) \leq \alpha \}.
    \end{align}
\end{Definition}
\subsection{Lower semicontinuity}
\begin{Definition}[Lower limit]
    For \(f: \mathbb{E}\to \mathbb{R}\), and \( \bar x \in \mathbb{E} \), 
    we define \begin{align}
        \liminf_{x\to \bar x} f(x) &= \sup_{\delta > 0} \inf_{x \in B(\bar x, \delta)} f(x)\\
        &= \inf\{\alpha \in \overline{\mathbb{R}} : \exists x^k \to x \text{ s.t. }
        f(x^k) \to \alpha 
        \}.
    \end{align}
    as the lower limit of \( f \) at \( \bar x \).
\end{Definition}

\begin{Definition}[Continuity notions]
    Let \( f:\mathbb{E}\to \overline{\mathbb{R}} \) 
    and \( \bar x \in \mathbb{E} \), then \( f \) is 
    said to be lower semicontinuous at \( \bar x \) if \( f(\bar x) \leq \liminf_{x\to \bar x} f(x) \).
\end{Definition}
Alternatively, \( f \) is lsc at \( \bar x \) iff there does not exists 
a sequence \( \{x^k\} \to \bar x \) such that \( f(x^k) < f(\bar x) \).

\begin{Definition}[Closure of function]
    For \( f:\mathbb{E}\to \overline{\mathbb{R}} \),
    the closure of \( f \) is the function \( \text{cl} f:\mathbb{E}\to \overline{\mathbb{R}} \) defined by
    \begin{align}
    \text{cl} f(\bar x) = \liminf_{x\to \bar x} f(x).
    \end{align}
\end{Definition}
As \( \liminf_{x\to \bar x} f(x) \leq f(\bar x) 
 \), we always have \( \text{cl} f \leq f \).
 We also have \(\text{epi}(\text{cl} f) = \text{cl}(\text{epi} f)\).
 That said, epigraph and closure commute with each other.

 \begin{Proposition}[Characterization of lsc]
     Let \( f:\mathbb{E}\to \overline{\mathbb{R}} \). Then TFAE:
     \begin{enumerate}
         \item \( f \) is lsc on \( \mathbb{E} \)
         \item \( \text{epi}(f) \) is closed
         \item \( \text{lev}_{\leq\alpha} f \) is closed for all \( \alpha \in \mathbb{R} \)
     \end{enumerate}
 \end{Proposition}

\begin{Corollary}[lsc of indicator function]
    Let \( C \subseteq \mathbb{E} \), its indicator function \( \delta_C \) is 
    proper, lsc iff \( C \) is nonempty and closed.
\end{Corollary}

 \subsection{Optimization problems}
 We always have 
 \begin{align}
    \inf_C f = - \sup_C (-f), \; \sup_C f = - \inf_C (-f).
 \end{align}

\begin{Definition}[Minimizer]
    Let \( f:\mathbb{E}\to \mathbb{R}\cup\{\infty\} \) and \( C \subseteq \mathbb{E} \).
    Then \( \bar x \in C \cap \text{dom}(f) \) is called 
    \begin{enumerate}
        \item local minimizer of \( f \) over \( C \) if 
        \begin{align}
            \exists \varepsilon > 0: f(\bar x) \leq f(x), \forall x \in C \cap B(\bar x, \varepsilon).
        \end{align}
        \item global minimizer of \( f \) over \( C \) if
        \begin{align}
            f(\bar x) \leq f(x), \forall x \in C.
        \end{align}
    \end{enumerate}
\end{Definition}

When minimizing a proper function \( f: \mathbb{E}\to \mathbb{R} \cup \{\infty\}\),
we have:
\begin{align}
    \arg\min_\mathbb{E} f \neq \emptyset \implies \inf_\mathbb{E} f \in \mathbb{R}.
\end{align}
but not the reverse.

\begin{Theorem}[Existence of minima]
    \label{thm:exi_min}
    Let \( f:\mathbb{E}\to \mathbb{R}\cup\{\infty\} \) be proper, 
    lsc, and level bounded. Then \( \arg\min f \neq \emptyset \)
    and \( \min f = \inf f \in \mathbb{R} \).
\end{Theorem}

\clearpage
\section{Convex Sets}
\begin{Definition}[Convex sets]
    A set \( C \subseteq \mathbb{E} \) is called convex if for all \( x, y \in C \) and \( \lambda \in [0, 1] \), we have
    \begin{align}
    \lambda x + (1-\lambda) y \in C.
    \end{align}
\end{Definition}
Example of convex sets includes:
subspace of \( \mathbb{E} \), Minkowski sum of convex sets, hyperplanes, 
affine sets, intersection of convex sets, linear image and 
preimage of convex sets, half spaces, interval (closed, open, half-open).

\subsection{Projection on convex sets}
\begin{Definition}[Projection on a set]
    let \( S\subset \mathbb{E} \) be nonempty and 
    \( x\in \mathbb{E} \), then we define the projection 
    of \( x \) on \( S \) by
    \begin{align}
        P_S(x) = \arg\min_{y\in S} \|x-y\| \subset S
    \end{align}
\end{Definition}

Now we give sufficient conditions for this subset to be nonempty and unique.
\begin{Lemma}
    Let \( x \in \mathbb{E} \) and \( S \subseteq \mathbb{E} \) be
    nonempty. Then the following hold:
    \begin{enumerate}
        \item If \( S \) is closed, then \(P_S(x) \neq \emptyset \).
        \item If \( S \) is convex, then \( P_S(x) \) has at most one element.
    \end{enumerate}
\end{Lemma}

\begin{Corollary}[Projection on closed, convex set]
    Let \( C \subset \mathbb{E} \) be nonempty, closed, and convex. 
    Then \( P_C(x) \) is a mapping from \( \mathbb{E} \) to \( C \)
    with \( x=P_C(x) \) iff \( x \in C \).
    That said, projection onto a closed convex set is unique.
\end{Corollary}

The following theorem gives an important characterization of the projection on a 
closed convex set in terms of variational inequality.

\begin{Theorem}[Projection Theorem]
    Let \( C \subseteq \mathbb{E} \) be nonempty, closed, and convex,
    and let \( x \in \mathbb{E} \). Then \(\bar v = P_C(x)\) iff
    \begin{align}
    \langle x - \bar v, v - \bar v \rangle \leq 0, \quad \forall v \in C.
    \end{align}
\end{Theorem}

\subsection{Separation Theorem}
\begin{Theorem}[Separation Theorem]
    Let \( C \subset \mathbb{E}\) be nonempty, closed, and convex,
    and let \( x\not\in C \), then there exists \( s \in \mathbb{E}\setminus \{0\}\) with 
    \begin{align}
        \langle s, x \rangle > \sup_{v \in C} \langle s, v \rangle.
    \end{align}
\end{Theorem}
Substituting \(s\) with \(-s\), we obtain 
\begin{align}
    \langle s, x \rangle < \inf_{v \in C} \langle s, v \rangle.
\end{align}


To see the separation, define \( \gamma = \frac12 (\langle s,x\rangle + \sup_{y\in C}\langle s,y\rangle)\).
Then, \( x\in \{
    z:\langle s, z\rangle > \gamma 
\} \) and \( C 
\subseteq \{ z : \langle s, z\rangle < \gamma \} \),
i.e., \(x\) and \( C \) 
lie in two distinct half-spaces defined by the hyperplane 
\begin{align}
    H = \{ z : \langle s, z\rangle = \gamma \}.
\end{align}


\clearpage
\section{Convex Functions}

\subsection{Definition and Examples}

\begin{Definition}[Convex functions]
    A function \( f:\mathbb{E}\to \overline{\mathbb{R}} \) is called
    convex if \(\text{epi}(f)\) is a convex set.
\end{Definition}
We can substitute the the epigraph with strict epigraph.

Convex function also have convex level sets.

Examples of convex functions: scalar-valued linear functions 
\(\mathcal{L}(\mathbb{E}_1, \mathbb{R})\) is convex; indicator of 
function of a convex set, any norm on \( \mathbb{E} \), the scalar functions
\(x\mapsto x^2\), \(x \mapsto \exp(x)\).

\begin{Proposition}[Domain of a convex function]
    The domain of a convex function is convex.
\end{Proposition}
Using the projection linear mapping \(L: (x,\alpha) \in \mathbb{E} \times \mathbb{R} 
\mapsto x\in \mathbb{E}\),
we have \(\text{dom}(f) = L(\text{epi}(f))\).
We know that linear map preserve convexity.

\begin{Proposition}[Characterizing convexity]
    A function \( f:\mathbb{E}\to \mathbb{R}\cup \{\infty\} \) is 
    convex iff for all \( x, y \in \mathbb{E}\) and \( \lambda \in (0, 1) \), we have
    \begin{align}
        f(\lambda x + (1-\lambda) y) \leq \lambda f(x) + (1-\lambda) f(y).
        \label{eq:convexity}
    \end{align}
\end{Proposition}
\marginnote{
\begin{Corollary}
    Let \( f:\mathbb{E}\to \mathbb{R}\cup \{\infty\} \), then TFAE:
    \begin{enumerate}
        \item \( f \) is convex
        \item \( f \) is convex on its domain
    \end{enumerate}
\end{Corollary}
We note that 
convexity on the domain is equivalent to convexity on the whole space 
because the inequality trivially holds when \( x \) or \( y \) is outside the domain.
}
We can extend the above characterization.

\begin{Definition}[p-dimensional simplex]
    The p-dimensional simplex is the set
    \begin{align}
        \Delta_p = \{ \lambda \in \mathbb{R}^p : \lambda_i \ge 0, \sum_{i=1}^p \lambda_i = 1 \}.
    \end{align}
\end{Definition}

\begin{Corollary}[Jensen's Inequality]
    Let \( f:\mathbb{E}\to \mathbb{R}\cup \{\infty\} \) be a convex function, and let \( x_1, \ldots, x_n \in \text{dom}(f) \) with \( \lambda_1, \ldots, \lambda_n \geq 0 \) such that \( \sum_{i=1}^n \lambda_i = 1 \). Then
    \begin{align}
        f\left(\sum_{i=1}^n \lambda_i x_i\right) \leq \sum_{i=1}^n \lambda_i f(x_i).
    \end{align}
\end{Corollary}

\begin{Definition}[Convexity on a set]
    A function \( f:\mathbb{E}\to \overline{\mathbb{R}} \) is convex on a nonempty convex set 
    \( C \subseteq \text{dom}(f) \)  if 
    the inequality \eqref{eq:convexity} holds for every \( x, y \in C \).
\end{Definition}

We mainly interest in proper convex (sometimes lsc) functions \(f: \mathbb{E}\to \mathbb{R}\cup \{\infty\} \).
We take the following notation:
\begin{itemize}
    \item \( \Gamma = \Gamma(\mathbb{E}) = \{f: \mathbb{E} \to \mathbb{R} \cup \{\infty\} : f \text{ is proper and convex}\}
    \) 
    \item \( \Gamma_0 = \Gamma_0(\mathbb{E}) = \{f: \mathbb{E} \to \mathbb{R} \cup \{\infty\} : f \text{ is proper, convex, and lsc}\}
    \) 
\end{itemize}
We then introduce stronger notions of convexity.
\begin{Definition}[Strict/strong convexity]
    Let \( f\in \Gamma \) and \( C \subset \text{dom}(f) \) a convex set, then \( f \) is said to be 
    \begin{itemize}
        \item strictly convex if for all \( x, y \in \mathbb{E} \) with \( x \neq y \) and all \( \lambda \in (0, 1) \), we have
    \begin{align}
        f(\lambda x + (1-\lambda) y) < \lambda f(x) + (1-\lambda) f(y).
    \end{align}
        \item strongly convex with parameter \( m > 0 \) if for all \( x, y \in \mathbb{E} \) and all \( \lambda \in (0, 1) \), we have
    \begin{align}
        f(\lambda x + (1-\lambda) y) \leq \lambda f(x) + (1-\lambda) f(y) - \frac{m}{2} \lambda (1-\lambda) \|x - y\|^2.
    \end{align}
    where \( m \) is called the modulus of strong convexity on \( C \)
    \end{itemize}
\end{Definition}
\marginnote{Note that strong convexity implies strict convexity, but not the reverse.}

\begin{Proposition}[Characterization of strong convexity]
    Let \( f \in \Gamma \) and \( C \subset \text{dom}(f) \). Then 
    \( f \) is strongly convex with parameter \( m > 0 \) on \( C \) iff the function \( x \mapsto f(x) - \frac{m}{2} \|x\|^2 \) is convex on \( C \).
\end{Proposition}

\subsection{Functional operations preserving convexity}
\begin{Proposition}[Positive combinations of convex functions]
    For \( p \in \mathbb{N} \), let \( f_1, \ldots, f_p \) be convex (and lsc) and 
    \(\alpha_i\geq 0\) for \(i=1,\ldots, p\). 
    Then 
    \begin{align}
    f = \sum_{i=1}^p \alpha_i f_i
    \end{align}
    is convex (and lsc).\\
    If, in addition, \(\bigcap_{i=1}^p \text{dom}(f_i) \neq \emptyset\), then \( f \) is also proper.
\end{Proposition}


\begin{Proposition}[Pointwise supremum of convex functions]
    For an aribitrary index set \( I\), let \( f_i \) be convex (lsc) for all \( i \in I \). 
    Then the function 
    \begin{align}
        f(x) = \sup_{i \in I} f_i(x)\; \forall x \in \mathbb{E}
    \end{align}
    is convex (lsc).
\end{Proposition}
\marginnote{Index does not have to be finite.}

\begin{Definition}[Affine functions]
    A function \( f:\mathbb{E}_1 \to \mathbb{E}_2 \) is called affine if 
    \begin{align}
    f(\lambda x + (1-\lambda) y) = \lambda f(x) + (1-\lambda) f(y), \forall x, y \in \mathbb{E}_1, \lambda \in [0, 1].
    \end{align}
\end{Definition}
\marginnote{If the coefficients need not sum to 1 then it's linear.
Linear function is affine and affine function is linear if \( f(0)=0 \).
}
Every affine function can have a unique representation 
\(f = L + v\) where \( L \in \mathcal{L}(\mathbb{E}_1, \mathbb{E}_2)
\) is linear and \( v \in \mathbb{E}_2 \).

\begin{Proposition}[Pre-composition of with an affine mapping]
    Let \(H: \mathbb{E}_1 \to \mathbb{E}_2\) be affine and \(g: \mathbb{E}_2 \to \mathbb{R}\cup \{\infty\}\) 
(lsc and ) convex. Then the function \( f = g \circ H \) is (lsc and) convex.
\end{Proposition}

This is not true in general, as the composition of two convex functions may not be convex.

\clearpage
\section{Minimization and convexity}
In an minimization problem:
\begin{align}
\min_{x \in C} f(x) \iff \min_{x \in \mathbb{E}} f(x) + \delta_C(x).
\end{align}
If \( f \) is convex and \( C \) is convex, 
then we call this problem a convex minimization/optimization problem.

\subsection{General existence results}
\begin{Definition}[Coercivity and supercoercivity]
    Let \( f:\mathbb{E}\to \overline{\mathbb{R}} \). Then \( f \) is called 
    \begin{itemize}
        \item coercive if \( \lim_{\|x\|\to \infty} f(x) = \infty \), which is sometimes called 0-coercive.
        \item supercoercive if \( \lim_{\|x\|\to \infty} \frac{f(x)}{\|x\|} = \infty \), which is sometimes called 1-coercive.
    \end{itemize}
\end{Definition}

\begin{Lemma}[Level-boundedness and coercivity]
    A function \( f:\mathbb{E}\to \overline{\mathbb{R}} \) is coercive iff \( f \) is level-bounded.
\end{Lemma}

We now extent \Cref{thm:exi_min} to cover the problem above.
\begin{Corollary}[Existence of minimizer]
    Let \( f:\mathbb{E}\to \mathbb{R}\cup \{\infty\} \) be lsc and let 
    \( C \subseteq \mathbb{E} \) be closed such that \( \text{dom}(f) \cap C \neq \emptyset \)
    and suppose one of the following holds:
    \begin{enumerate}
        \item \( f \) is coercive
        \item \( C \) is bounded
    \end{enumerate}
    Then \( f \) has a minimizer on \( C \).
\end{Corollary}
Proof: Consider the function \( g = f + \delta_C \).
Then it holds that \(\text{lev}_{\leq \alpha} g = \text{lev}_{\leq \alpha} f \cap C\) for all \( \alpha \in \mathbb{R} \).

Hence, under either assumption, \( g \) is closed and bounded level sets.
Hence, it is lsc and level-bounded, Then the results follows from \Cref{thm:exi_min}.
\hfill\qedsymbol




\clearpage
\section{Differential Theory}
\subsection{Subdifferential}
\begin{Definition}[Subdifferential]
    Let \( f:\mathbb{E}\mapsto \overline{\mathbb{R}} \), and \( \bar x \in \mathbb{E} \).
    Then \( g\in \mathbb{E} \) is call a subgradient of \( f \) at \( \bar x \) 
    if 
    \begin{align}
    f(x) \geq f(\bar x) + \langle g, x - \bar x \rangle, \quad \forall x \in \mathbb{E}.
    \end{align}
    The set of all subgradients of \( f \) at \( \bar x \):
    \begin{align}
    \partial f(\bar x) = \{ g \in \mathbb{E} : f(x) \geq f(\bar x) + \langle g, x - \bar x \rangle, \forall x \in \mathbb{E} \}
    \end{align}
    is called the subdifferential of \( f \) at \( \bar x \).
\end{Definition}
Subdifferential contains exactly the slopes of all affine minorants
of \( f \) at \( \bar x \).

\begin{Proposition}[Elementary properties of the subdifferential]
    Let \( f:\mathbb{E}\mapsto \overline{\mathbb{R}} \). Then the following hold:
    \begin{enumerate}
        \item \( \partial f(\bar x) \) is closed and convex for all \( \bar x \in \text{dom}(f) \).
        \item If \( f \) is proper, then \( \partial f(\bar x) = \emptyset \) for all \( \bar x \notin \text{dom}(f) \).
        \item \( 0 \in \partial f(\bar x) \) iff \(\bar x \in \arg\min f\).
        \item If \( f \) is convex, then \( \partial f(\bar x) = \{ v \in \mathbb{E} :
        (v, -1) \in N_{\text{epi}(f)}(\bar x, f(\bar x)) \} \).
    \end{enumerate}
\end{Proposition}

\begin{Theorem}[Subdifferential and conjugate function]
    Let \( f:\mathbb{E}\mapsto \overline{\mathbb{R}} \). Then 
    the following are equivalent:
    \begin{enumerate}
        \item \( y \in \partial f (x) \).
        \item \( x \in \arg\min_z\{\langle z, y \rangle - f(z)\} \)
        \item \( f(x) + f^*(y) = \langle x, y \rangle \) \textbf{[Fenchel-Young equality]}
    \end{enumerate}
    If \( f \in \Gamma_0\), these are also equivalent to
    \begin{enumerate}
        \setcounter{enumi}{3}
        \item \( x \in \partial f^*(y) \)
        \item \( y \in \arg\min_z\{\langle x, z \rangle - f^*(z)\} \)
    \end{enumerate}
    Note that (1) and (4) are equivalent to \( \partial f^* \) 
    is an inverse of \( \partial f \).
\end{Theorem}
Proof: Notice that 
\begin{align}
    y \in \partial f(x) &\iff f(z) \geq f(x) + \langle y, z - x \rangle, \forall z \in \mathbb{E} \\
    &\iff \langle z, y \rangle - f(z) \le \langle x, y \rangle - f(x), \forall z \in \mathbb{E} \\
    &\iff \langle x, y \rangle - f(x) \ge \sup_z \{\langle z, y \rangle - f(z)\} = f^*(y) \\
    \intertext{Combine with Fenchel-Young inequality \( f(x) + f^*(y) \ge \langle x, y \rangle \), we have}
    &\iff f(x) + f^*(y) = \langle x, y \rangle
\end{align}
Apply the same argument to \( f^* \) and using the fact that \( f^{**} = f \) for \( f \in \Gamma_0 \), we 
complete the proof.
\hfill\qedsymbol

\subsection{Directional Derivative of a convex function}
\begin{Definition}[Directional derivative]
    Let \( f:\mathbb{E}\mapsto \overline{\mathbb{R}} \) be proper, and \( \bar x \in \text{dom}(f) \). 
    We say that \( f \) is directionally differentiable at \( \bar x \) 
    in the direction \( d \in \mathbb{E} \)
    if the limit
    \begin{align}
    f'(\bar x; d) = \lim_{t\downarrow 0} \frac{f(\bar x + td) - f(\bar x)}{t}.
    \end{align}
    exists in an extended real-valued sense (so it can be \( +\infty \)).
    In this case, we call \( f'(\bar x; d) \) the directional derivative of 
    \( f \) at \( \bar x \) in the direction \( d \).
\end{Definition}
\marginnote{
Remark: if \( f \) is differentiable, then \( f'(\bar x; d) = \langle \nabla f(\bar x), d \rangle \).
}

\begin{Definition}[Positive Homogeneity]
    We call a function \( f:\mathbb{E}\mapsto \mathbb{R}\cup \{+\infty\} \) 
    positively homogeneous if \( f(\lambda x) = \lambda f(x) \) for all \( x \in \mathbb{E} \) and \( \lambda \ge 0 \).
\end{Definition}
\begin{Proposition}
    A function \( f:\mathbb{E}\mapsto \mathbb{R}\cup \{+\infty\} \) is
    convex and positively homogeneous if and only if
    \begin{align}
        h(\lambda x + \mu y) \le \lambda h(x) + \mu h(y), \quad \forall x, y \in \mathbb{E}, \lambda, \mu \ge 0.
    \end{align}
    a property called sublinearity. Sublinearity
    is also equivalent to subadditivity \( h(x + y) \le h(x) + h(y) \) and positive homogeneity.
\end{Proposition}
That said, convexity and positive homogeneity give subadditivity:
\begin{align}
    f(\lambda x+ (1-\lambda) y) \le \lambda f(x) + (1-\lambda) f(y), \quad \forall x, y \in \mathbb{E}, \lambda \in [0, 1].\\
    \intertext{Set \( \lambda = \frac{1}{2} \), we have}
    f\left(\frac{x+y}{2}\right) \le \frac{f(x) + f(y)}{2} \implies f(x+y) \le f(x) + f(y).
\end{align}
\hfill \qedsymbol

Observe that any positive homogeneous function has \( h(0) = 0 \) and
any sublinear function satisfies \( -h(x) \le h(0) + h(-x) \le h(x) \).

\begin{Proposition}[Directional derivative of a convex function]
    Let \( f\in \Gamma_0 \), \( x\in \text{dom}(f) \), and \( d \in \mathbb{E} \).
    Define:
    \begin{align}
        t\in \mathbb{R} \setminus \{0\} \mapsto q(t) = \frac{f(x + td) - f(x)}{t}.
    \end{align}
    then TFAE:
    \begin{enumerate}
            \item \( q(-t) \leq q(-s) \leq q(s) \leq q(t) \) for all \( 0<s<t \).
            \item We have \( f'(x; d) = \inf_{t>0} q(t) \). In particular, \( f'(x; d) \) always exists in an extended real-valued sense.
            \item \( \text{dom}(f') = \mathbb{R}_+(\text{dom}(f) - x) \).
            \item If \( x \in \text{int}(\text{dom}(f)) \), then \( f'(x; \cdot) \) is finite-valued, positively homogeneous,
            and convex, (so sublinear).
        \end{enumerate}
\end{Proposition}
\marginnote{For notation: \\
\( \mathbb{R}_+(S) := \{ \alpha s \mid s \in S, \alpha \ge 0 \} \).
\( d \in \mathbb{R}_+(S) \iff \exists\alpha \ge 0, s \in S : d = \alpha s \).
}
Proof:
(2): Since infima always exists in \( \overline{\mathbb{R}} \), \( \inf_{t>0}q(t) \)
exists. By (1), \( q(t) \) monotonically decreasing as \( t \downarrow 0 \),
so we have 
\begin{align}
    \inf_{t>0} q(t) = \lim_{t\downarrow 0} q(t) = f'(x; d).
\end{align}
(3): Follows from (2) we have:
\begin{align}
    d \in \text{dom}(f') &\iff \exists  t > 0  : \frac{f(x+td) - f(x)}{t} < \infty.\\
    &\iff \exists  t > 0  : f(x + td) - f(x) < \infty.\\
    &\iff \exists  t > 0  : x + td \in \text{dom}(f), \text{ since \( f(x) < \infty \) for \( x \in \text{dom}(f) \)}.\\
    &\iff \exists  t > 0  : td \in \text{dom}(f) - x.\\
    &\iff d \in \mathbb{R}_+(\text{dom}(f) - x).
\end{align}
(4): Take \( 0<s<t \). If \( x \in \text{int}(\text{dom}(f)) \),
we can choose \( t>0 \) sufficiently small such that both \( x\pm td \in \text{int}(\text{dom}(f)) \).
Thus \( q(t) \) and \( q(-t) \) are finite. 
Hence \( +\infty > q(t) \geq q(s) \downarrow f'(x;d) \geq q(-t) > -\infty \).
As \( d \in \mathbb{E} \) is arbitrary, \( f'(x; \cdot) \) is finite-valued.

\( f'(x,\cdot) \) is also positively homogeneous as:
\begin{align}
    f'(x; \alpha d) &= \lim_{t\downarrow 0} \frac{f(x + t\alpha d) - f(x)}{t} \\
    &= \lim_{t\downarrow 0} \alpha \frac{f(x + t\alpha d) - f(x)}{t\alpha} \\
    &= \alpha f'(x; d).
\end{align}

Convexity of \( f'(x; \cdot) \):
for \( t>0, d,h \in \mathbb{E} \) and \( \lambda \in (0,1) \), 
by convexity of \( f \) we have
\begin{align}
    f(x + t(\lambda d + (1-\lambda) h)) -f(x) 
    &= f(\lambda (x + td) + (1-\lambda)(x + th)) - f(x) \\
    &\le 
    \lambda (f(x + td) - f(x)) + (1-\lambda)(f(x + th) - f(x)).
\end{align}
which implies
\begin{align}
    f'(x; \lambda d + (1-\lambda) h) &= \lim_{t\downarrow 0} \frac{f(x + t(\lambda d + (1-\lambda) h)) - f(x)}{t} \\
    &\le \lim_{t\downarrow 0} \lambda \frac{f(x + td) - f(x)}{t} + (1-\lambda)\frac{f(x + th) - f(x)}{t} \\
    &= \lambda f'(x; d) + (1-\lambda) f'(x; h).
\end{align}
\hfill\qedsymbol

\begin{Definition}[Linearity space]
    For a sublinear function \( h:\mathbb{E}\mapsto \mathbb{R}\cup \{+\infty\} \), we call the set
    \begin{align}
        \text{lin}(h) = \{ d \in \mathbb{E} :- h(d) = h(-d)  \}
    \end{align}
    the linearity space of \( h \).
\end{Definition}

\begin{Lemma}[Linearity space of a sublinear function]
    Let \( h:\mathbb{E}\mapsto \mathbb{R}\cup \{+\infty\} \) be sublinear. Then 
    the following holds:
    \begin{enumerate}
        \item The linearity space \( \text{lin}(h) \) is a subspace on which \( h \) is linear.
        \item For \( \bar x \in \text{int}(\text{dom}(h)) \) and \( q=h'(\bar x, \cdot) \), we have:
        \begin{enumerate}
            \item \(q(\lambda \bar x) = \lambda h(\bar x)\) 
            \item \( q\leq h \)
            \item \(\text{lin} q \supset \text{lin} h + \text{span}\{\bar x\}\)
        \end{enumerate}
    \end{enumerate}
\end{Lemma}

\begin{Theorem}[Subdifferential vs. directional derivative]
    Let \( f\in \Gamma \), then the following holds:
    \begin{enumerate}
        \item For all \( \bar x \in \mathbb{E}\), we have 
        \(\partial f(\bar x) = \{s \in \mathbb{E}:\langle s,d\rangle \leq f'(\bar x , d), d \in \mathbb{E}\}
        \).
        \item For all \( \bar x \in \text{int}(\text{dom}(f)) \), we have
        \begin{align}
            f'(\bar x; d) = \sup_{s \in \partial f(\bar x)} \langle s, d \rangle, \quad \forall d \in \mathbb{E}.
        \end{align}
    \end{enumerate}
    In particular, \( \partial f(\bar x) \) is nonempty.
\end{Theorem}
\subsection{Continuity properties}
The subdifferential operator of a closed, proper, and convex function has a closed graph.
\begin{Corollary}[Outer semicontinuity of the subdifferential]
    Let \( f\in \Gamma_0 \) and suppose \( \{x_k\} \to x \) and 
    \( \{y_k \in \partial f(x_k)\} \to y \). Then \( y \in \partial f(x) \), i.e. the graph of \( \partial f \) is closed.
\end{Corollary}
Proof: By Fenchel-Young equality, we have \( f(x_k) + f^*(y_k) = \langle x_k, y_k \rangle \) for all \( k \). 
On the other hand, by lsc of \( f \) and \( f^* \), we have
\begin{align}
    f(x) + f^*(y) \leq \liminf_{k\to\infty} f(x_k) + \liminf_{k\to\infty} f^*(y_k) \leq \liminf_{k\to\infty} (f(x_k) + f^*(y_k)) = \liminf_{k\to\infty} \langle x_k, y_k \rangle = \langle x, y \rangle.
\end{align}
\( f(x) + f^*(y) \leq \langle x,y\rangle \). Together with Fenchel-Young inequality \( f(x) + f^*(y) \geq \langle x, y \rangle \), we have \( f(x) + f^*(y) = \langle x, y \rangle \).
This implies \( y \in \partial f(x) \).
\hfill\qedsymbol

A proper convex function is locally Lipschitz continuous on the interior of its domain.
\begin{Definition}[Locally Lipschitz continuous]
    A function \( f:\mathbb{E}\mapsto \mathbb{R}\cup \{+\infty\} \) is locally Lipschitz continuous at \( \bar x \in \text{dom}(f) \) 
    if \(\exists L_{\bar x}, \delta_{\bar x} > 0\) such that 
    \begin{align}
    |f(x) - f(y)| \leq L_{\bar x} \|x-y\|, \quad \forall x, y \in B(\bar x, \delta_{\bar x}).
    \end{align}
\end{Definition}
Remark: every function that is locally Lipschitz at a point is also continuous at that point.

If a function is locally Lipschitz at every point of a compact set, then it is globally Lipschitz on that set.

\begin{Theorem}[Local Lipschitz continuity of a proper convex function]
    Let \( f\in \Gamma \) and \( \bar x \in \text{int}(\text{dom}(f)) \). Then \( f \) is locally Lipschitz continuous at \( \bar x \).
\end{Theorem}

\begin{Definition}[Domain of continuity]
    For a function \( f:\mathbb{E}\mapsto \mathbb{R}\cup \{+\infty\} \), we call the set
    \begin{align}
    \text{cont}(f) = \{ x \in \mathbb{E} : f \text{ is continuous at } x \}
    \end{align}
    the domain of continuity of \( f \).
\end{Definition}

\begin{Corollary}
    Let \( f\in \Gamma \). Then \( \text{int}(\text{dom}(f)) = \text{cont}(f) \).
\end{Corollary}

\begin{Theorem}[Boundedness properties of the subdifferential]
    For \( f\in \Gamma \), the following holds:
    \begin{enumerate}
        \item Let \(X \subset \text{int}(\text{dom}(f))\) be nonempty, open, and convex. Then \( f \)
        is Lipschitz continuous with modulus \( L >0 \) on \( X \) if and only if \( \|v\| \leq L \) for all 
        \( v \in \partial f(x) \) and \( x \in X \).
        \item \(\partial f\) maps bounded sets which are compactly contained in \( \text{int}(\text{dom}(f)) \) to bounded sets.
    \end{enumerate}
\end{Theorem}

\begin{Corollary}
    Let \( f\in \Gamma \) and \( \bar x \in \text{int}(\text{dom}(f)) \). Then \( \partial f(\bar x) \) is nonempty, compact, and convex.
\end{Corollary}
Proof: By elementary properties of the subdifferential, \( \partial f(\bar x) \) is closed and convex since 
\( \bar x \in \text{int}(\text{dom}(f)) \). By the boundedness properties of the subdifferential, 
since \( \{\bar x\} \) is compactly contained in \( \text{int}(\text{dom}(f)) \), 
\( \partial f(\bar x) \) is bounded. Hence \( \partial f(\bar x) \) is compact.
\hfill\qedsymbol


\subsection{Differentiable convex function}
Differentiability of a convex function is defined on an open set.
\begin{Theorem}
    Let \( f\in \Gamma \) and \( \bar x \in \text{int}(\text{dom}(f)) \). Then
    \( \partial f(\bar x) \) is a singleton iff \( f \) is differentiable at \( \bar x \). 
    In this case, \( \partial f(\bar x) = \{\nabla f(\bar x)\} \).
\end{Theorem}
Differentiability on an open set implies continuous differentiability on that set.
\begin{Theorem}
    Let \( f\in \Gamma \) and \( \bar x \in \text{int}(\text{dom}(f)) \). Then \( f \) is continuously
    differentiable at \( \bar x \) iff \( \partial f \) is a singleton for all \( x \in \text{int}(\text{dom}(f)) \) and the mapping \( x \mapsto \partial f(x) \) is continuous at \( \bar x \). 
\end{Theorem}

\begin{Corollary}
    Let \( f: \mathbb{E}\to \mathbb{R} \cup \{+\infty\} \) be convex. Then \( f \) 
    if differentiable iff \( f \) is continuously differentiable.
\end{Corollary}

\subsubsection{Second-order characterizations}
\begin{Theorem}
    Let \( f: \mathbb{E}\to \mathbb{R} \cup \{+\infty\} \) be twice differentiable 
    on the open convex set \( \Omega \subset \text{int}(\text{dom}(f)) \).
    Then the following hold:
    \begin{enumerate}
        \item \( f \) is convex on \( \Omega \) iff \( \nabla^2 f(x) \succeq 0 \) for all \( x \in \Omega \).
        \item If \( \nabla^2 f(x) \) is positive definite for all \( x \in \Omega \), then \( f \) is strictly convex on \( \Omega \).
        \item \( f \) is \( \sigma \) strongly convex on \( \Omega \) iff for all \( x \in \Omega \), 
        the smallest eigenvalue of \( \nabla^2 f(x) \) is bounded below by \( \sigma \), i.e. \( \nabla^2 f(x) \succeq \sigma I \).
    \end{enumerate}
\end{Theorem}
\begin{Corollary}
    Let \( f: \mathbb{R}^n \to \mathbb{R}\) be defined by 
    \begin{align}
        f(x) = \frac{1}{2} x^T A x + b^T x + c,
    \end{align}
    with \( A \in \mathbb{R}^{n\times n} \) symmetric, \( b \in \mathbb{R}^n \), and \( c \in \mathbb{R} \). Then \( f \) 
    is convex iff \( A \succeq 0 \), and \( f \) is \( \sigma \)-strongly convex iff \( f \) is 
    strictly convex iff \( A \) is positive definite.
\end{Corollary}

\subsection{Subdifferential calculus}
\begin{Theorem}
    For all \( i \in I={1,\ldots,m} \), let \( f_i \in \Gamma\) and \( x \in \bigcap_{i=1}^m \text{int}(\text{dom}(f_i)) \).
    Set \( f = \max_{i\in I} f_i \) and \( I(x) = \{ i \in I : f_i(x) = f(x) \} \). Then
    \begin{align}
        \partial f(x) = \text{conv}\left( \bigcup_{i\in I(x)} \partial f_i(x) \right).
    \end{align}
\end{Theorem}

\begin{Corollary}
    For all \( i \in I={1,\ldots,m} \), let \( f_i \in \Gamma\) be 
    differentiable at \( x \in \bigcap_{i=1}^m \text{int}(\text{dom}(f_i)) \).
    Set \( f = \max_{i\in I} f_i \) and \(I(x) = \{ i \in I : f_i(x) = f(x) \} \). Then
    \begin{align}
        \partial f(x) = \text{conv}\left( \bigcup_{i\in I(x)} \{\nabla f_i(x)\} \right).
    \end{align}
\end{Corollary}

\begin{Theorem}[Fenchel-Rockafellar duality]
    Let \( f: \mathbb{E}_1 \to \mathbb{R}\cup \{+\infty\} \) and \( g: \mathbb{E}_2 \to \mathbb{R}\cup \{+\infty\} \)
and \( L \in \mathcal{L}(\mathbb{E}_1, \mathbb{E}_2) \). Then the followings hold:
    \begin{enumerate}
        \item (Weak duality)
        We have 
        \begin{align}
            p := \inf_{x\in \mathbb{E}_1}\{f(x) + g(L(x))\} \geq \sup_{y\in \mathbb{E}_2}\{-f^*(L^*(y)) - g^*(-y)\} =: d.
        \end{align}
        \item (Strong duality) If \( f,g \) are convex and the qualification condition \( 0 \in \text{int}(\text{dom}(g) - L(\text{dom}(f))) \) holds, 
        then equality holdes, i.e, \( p = d \), and the supremum is attained in finite.
        \item (Primal-dual recovery) If \( f\in \Gamma_0(\mathbb{E}_1) \), \( g \in \Gamma_0(\mathbb{E}_2) \),
        the following are equivalent for \( \bar x \in \mathbb{E}_1 \) and \( \bar y \in \mathbb{E}_2 \):
        \begin{enumerate}
            \item \( p=d \), \( \bar x \in \arg\min\{f(x) + g(L(x))\}\), and \( \bar y \in \arg\max\{-f^*(L^*(y)) - g^*(-y)\} \).
            \item \( L^*(\bar y)\in \partial f (\bar x) \), and \( -\bar y \in \partial g(L(\bar x)) \).
            \item \(\bar x \in \partial f^*(L^*(\bar y)) \), and \( L(\bar x) \in \partial g^*(-\bar y) \).
    \end{enumerate}
    \end{enumerate}
\end{Theorem}
We note that the qualification condition is satisfied under a stronger condition 
\( \text{int}(\text{dom}(g)) \cap L(\text{dom}(f)) \neq \emptyset \), which is often easier to verify in practice.

We call \( \inf_{x \in \mathbb{E}_1} \{f(x) + g(L(x))\} \) the primal problem, and \( \sup_{y \in \mathbb{E}_2} \{-f^*(L^*(y)) - g^*(-y)\} \) the dual problem.
Part (2), (3) of the theorem say under the qualification condition, 
solving the primal problem is equivalent to solving the dual problem, and we can get the 
\(\arg\min, \arg\max\) from (3).

Fenchel duality gives us the sum rule for subdifferential:
\begin{Theorem}[Generalized sum rule]
    For \( f: \mathbb{E} \to \mathbb{R}\cup \{+\infty\} \) and \( g: \mathbb{E} \to \mathbb{R}\cup \{+\infty\} \), 
    and \( L \in \mathcal{L}(\mathbb{E}, \mathbb{E}) \), the following hold:
    \begin{enumerate}
        \item \( \partial (f + g\circ L)(x) \supset \partial f(x) + L^*(\partial g(L(x))) \) for all \( x \in \mathbb{E} \).
        \item If \( f\in \Gamma_0(\mathbb{E}) \), \( g \in \Gamma_0(\mathbb{E}) \), and 
        the qualification condition \( 0 \in \text{int}(\text{dom}(g) - L(\text{dom}(f))) \) holds, 
        the equality holds, i.e. \( \partial (f + g\circ L)(x) = \partial f(x) + L^*(\partial g(L(x))) \) for all \( x \in \mathbb{E} \).
    \end{enumerate}
\end{Theorem}

\begin{Corollary}[Subdifferential sum rule]
    For \( f, g \in \Gamma\), then 
    \begin{align}
        \partial (f + g)(x) \supset \partial f(x) + \partial g(x), \quad \forall x \in \mathbb{E}.
    \end{align}
    If the qualification condition holds, equality holds, i.e. \( \partial (f + g)(x) = \partial f(x) + \partial g(x) \) for all \( x \in \mathbb{E} \).
\end{Corollary}

\begin{Corollary}[Subdifferential chain rule]
    Let \( g\in \Gamma(\mathbb{E}_2) \) and \( L \in \mathcal{L}(\mathbb{E}_1, \mathbb{E}_2) \).
    Then \( \partial (g\circ L)(x) \supset L^*(\partial g(L(x))) \). Under the 
    qualification condition, the equality holds, i.e. \( \partial (g\circ L)(x) = L^*(\partial g(L(x))) \) for all \( x \in \mathbb{E}_1 \).
\end{Corollary}

\begin{Theorem}[Subdifferential of infimal convolution]
    Let \( f, g \in \Gamma_0\) and \( x \text{dom}(f \# g)\). Then 
    \begin{align}
        \partial (f \# g)(x) = \partial f(\bar u) \cap \partial g(x - \bar u),
    \end{align}
    for all \( \bar u \in \arg\min_u \{f(u) + g(x-u)\} \).
\end{Theorem}

\begin{Corollary}[Derivative of Mereau envelop]
    Let \( f \in \Gamma_0 \) and \( \lambda>0 \). Then 
    \( e_\lambda f \) is continuously differentiable with \( \frac1\lambda \) Lipschitz gradient
    given by 
    \begin{align}
        \nabla e_\lambda f(x) = \frac{1}{\lambda}(I - \text{prox}_{\lambda f})(x).
    \end{align}
\end{Corollary}

\begin{Theorem}[Moreau decomposition]
    Let \( f \in \Gamma_0 \) and \( \lambda > 0 \). Then for all \( x \in \mathbb{E} \), we have
    \begin{align}
        x = \text{prox}_{\lambda f}(x) + \lambda \text{prox}_{\frac{1}{\lambda} f^*}\left(\frac{x}{\lambda}\right).
    \end{align}
\end{Theorem}

\begin{Proposition}
    For any proper, convex, closed function \( f:\mathbb{E}\to \overline{\mathbb{R}} \) and \(g: \mathbb{E} \to \mathbb{R} \) is 
    continuously differentiable, then
    \begin{align}
        \partial (f + g)(x) = \partial f(x) + \nabla g(x), \quad \forall x \in \mathbb{E}.
    \end{align}
\end{Proposition}

\begin{Corollary}[Primal-dual optimality condition]
    Consider the problems
    \begin{enumerate}
        \item[P.] \(\min_x h(Lx) + g(x)\)
        \item[D.] \(\max_y -g^*(-L^*(y))-h^*(y)\)
    \end{enumerate}
    where \( g: \mathbb{E} \to \overline{\mathbb{R}}\) and \( h: Y \to \overline{\mathbb{R}} \) 
    are proper, closed, and convex functions, and \( L \in \mathcal{L}(\mathbb{E}, Y) \).

    Suppose the optimal value of the primal and dual problem are equal, as implied by the qualification condition.
    Then, \( x \) is the minimizer of the primal problem and \( y \) is the maximizer of the dual problem iff
    \begin{align}
        \begin{bmatrix}
            0 \\ 0
        \end{bmatrix}
        \in \begin{bmatrix}
            0 & L^* \\
            -L & 0
        \end{bmatrix}
        \begin{bmatrix}
            x \\ y
        \end{bmatrix}
        +
        \partial g(x) \times \partial h^*(y).
    \end{align}
\end{Corollary}

% end of main text

\makeatletter
  \renewcommand{\section}{\@startsection{section}%
    {3}{0.8em}{-3ex \@plus -1ex \@minus -.2ex}%
    {1.5ex \@plus .2ex}
    {\hspace*{-5.5em}\fcolorbox{Periwinkle}{Periwinkle}{\parbox[c][1.0ex][b]{4em}{\phantom{space}}}
    \normalfont\Large\itshape\color{blue}}}
\makeatother

\bibliography{marginnotes}
\bibliographystyle{plainnat}

\end{document}
